%!TEX root = ./../../thesis-trier.tex

\IN{partial deterministic automaton}
\index{partial automaton|see{partial deterministic automaton}}
\index{PDFA|see{partial deterministic automaton}}
%\IN{state set in PDFAs}
%\IN{final state set in PDFAs}
\IN{state set!of a PDFA}
\IN{final state set!of a PDFA}
\IN{start state!of a PDFA}
\IN{transition function!of a PDFA}
\IN{input alphabet!of a PDFA}
Here, we generalize the notion of a complete and deterministic automaton.
A \emph{partial deterministic and finite
	automaton (PDFA)}
is a $5$-tuple $\automat{A}{Q}{\Sigma}{\delta}{q_0}{F}$
where, as for complete and deterministic automata, $Q$ is the \emph{finite state set},~$\Sigma$
the \emph{finite input alphabet}, $q_0$ a \emph{start state}, $F \subseteq Q$ a set of \emph{final states}, but where
 the \emph{transition function} is allowed to be a partial function $\delta \colon Q \times \Sigma \nrightarrow Q$.\footnote{We only consider PDFAs as language acceptors and hence always endow them with a start state and a set of final states.}
So, every complete and deterministic automaton is a partial automaton as well.\footnote{This is actually a rather unfortunate naming. Usually,
	one would expect that when adding modifiers to a concept
	this would narrow down the possibilities, which also allows the description of subtypes by combining modifiers,
	but here adding	the modifier ``partial'' gives a more general concept.
	For this reason, some authors, for example~\cite{DBLP:journals/tcs/CaiF22},
	deviate from the standard terminology by calling the most general model simply finite automaton.}
The modifier \emph{complete}
refers to the fact that the transition function is a total function, i.e., defined for each argument. 

\smallskip
\begin{center}
	\fbox{ 
		\begin{minipage}{28.5em}								
			When we just say ``automaton'' we always mean a complete and deterministic semiautomaton, i.e., 
			the automaton model
			previously considered with totally defined transition functions, and we explicitly
			say PDFA to denote a partial automaton with a start state and set of final states. In the context
			of this work, we only use PDFAs to describe constraint languages. Input automata
			to our problems are always assumed to be complete and deterministic.
		\end{minipage}
	}
\end{center}
\smallskip

\IN{accessible!PDFA}
\IN{accessible!state!in PDFA}
\IN{coaccessible state!in PDFA}
\IN{reachability order!PDFA}
\IN{accepted!language!by PDFA}
\IN{accepted!word!by PDFA}
\IN{PDFA!as language acceptor}
The PDFA $\mathcal A$ can be made complete by adding a non-final trap state
and by adding transitions to this new state for each~$q \in Q$ and $x \in \Sigma$
where $\delta(q, x)$ is undefined.
The notions accessible states, coaccessible states, accessible PDFAs, the reachability order on the states $\lereach$, accepted words,
the accepted language~$L(\mathcal A)$, the notation $\mathcal A_{q,S}$ for $q \in Q$ and $S \subseteq Q$, 
the extension of the partial transition function to words, i.e., 
$\delta \colon Q \times \Sigma^* \nrightarrow Q$, and the extension of~$\delta$ to subsets of states 
$\applytrans{\delta}{S}{u} = \{ \applytrans{\delta}{q}{u} \mid \applytrans{\delta}{q}{u} \text{ is defined and } q \in S \}$
all carry over to PDFAs without difficulty.\footnote{Martyugin~\cite{Martyugin14} defines the extension of the transition function to subsets differently, requiring that $\delta(q,u)$ be defined for every $q \in S$, and leaving $\delta(S,u)$ undefined otherwise. A similar approach for nondeterministic automata was used in~\cite[Sec.~8]{DBLP:conf/aaai/ArrighiFO023} in the context of \emph{conformant planning}. However, apart from these cases, this convention appears to be uncommon in the literature.}
We call~$S \subseteq Q$ an upper set if $\delta(S, x) \subseteq S$ for each $x \in \Sigma$.
We define $\mathcal A_{|S}$, $\mathcal A_{|\Gamma}$,
and $\mathcal A_{|\Gamma,S}$ by restricting the partial transition
function to inputs from $S$,$\Gamma$, and $S \times \Gamma$, respectively,
analogously to the corresponding constructions for the previous automaton model.
The product, power, and pair automata constructions also carry over and will be used in the following
for PDFAs.
A PDFA is called \emph{trim} if every state is accessible and coaccessible. Every non-empty regular language is accepted by some trim PDFA; the empty language $\emptyset$ is the only language that is not accepted by a trim PDFA.
\IN{trim}

\begin{proposition}
	\label{constrsync:prop:trim}
	Let $\automat{A}{Q}{\Sigma}{\delta}{q_0}{F}$
	be a PDFA. Then $\mathcal A$
	is trim if and only if~$[q_0]_{\eqreach}$
	is a minimum in the order induced by $\lereach$
	on the strongly connected components
	and $[q]_{\eqreach} \cap F \ne \emptyset$
	for every maximal strongly connected component $[q]_{\eqreach}$.	
\end{proposition}
\begin{proof}
	Every state is accessible if and only if $q_0 \lereach q$ for all $q \in Q$, i.e., $[q_0]_{\lereach}$
	is a minimum.
	Further, every state is coaccessible
	if and only if for each $q \in Q$ 
	we have $q \lereach p$ for some $p \in F$,
	which is fulfilled if every maximal strongly connected component contains a state from~$F$.
\end{proof}

Minimal PDFAs are characterized by adding the additional requirement that every state is coaccessible
to \cref{prelim:thm:char_minaut}.
The minimal PDFA differs from a minimal complete
automaton if and only if the minimal complete automaton has a non-final trap state,
which can be discarded in this case to obtain the minimal PDFA. 
Note that in a minimal complete automaton for the set of synchronizing words of a given complete and deterministic automaton
every state is coaccessible, so that  it makes no difference 
whether we consider complete or partial automata when investigating minimal deterministic automata for the 
for the set of synchronizing words.

\IN{exactly synchronizing word}
\IN{carefully synchronizing words}
\IN{careful synchronization of automaton}
\IN{exact synchronization of automaton}
Next, we introduce two generalizations of synchronization to PDFAs.
Let $u \in \Sigma^*$.
We call $u$ a \emph{carefully synchronizing word}~\cite{Martyugin14,Rystsov1980} for $\mathcal A$
if there exists $p \in Q$ such that $\delta(q, u) = p$
for all $q \in Q$, i.e., $|\delta(Q, u)| = 1$ and $\delta(q, u)$ is defined
for each $q \in Q$.
We call $u$ an \emph{exactly synchronizing word}~\cite{Shabana_2019,TraverCrutchfield2011}
if $|\delta(Q, u)| = 1$.
Further, $\mathcal A$ is said to be \emph{carefully synchronizing} (resp.~\emph{exactly synchronizing})
if there exists a carefully (resp.~exactly) synchronizing word for~$\mathcal A$.